\documentclass[a4paper,11pt]{article}

\usepackage[margin=2.5cm]{geometry}
\usepackage{listings}
\usepackage{hyperref}

\title{Practical Work 2: RPC File Transfer}
\author{Nguyen Trong Bach \\ Student ID: 23BI14057}
\date{\today}

\begin{document}

\maketitle

\section{Introduction}
This practical upgrades the TCP file transfer system from Practical Work 1 into an RPC-based system. Remote Procedure Call (RPC) allows a client to call functions on the server as if they were local. Python's XML-RPC module is used due to its simplicity and built-in support.

\section{RPC Service Design}
The RPC server exposes four methods:
\begin{itemize}
    \item \texttt{list\_files()} — lists available files
    \item \texttt{upload\_file(filename, data)} — uploads a file
    \item \texttt{download\_file(filename)} — returns file contents
    \item \texttt{delete\_file(filename)} — removes a file
\end{itemize}

The architecture is shown in Figure~1.

\section{System Organization}
The system is organized into:
\begin{enumerate}
    \item \textbf{server.py}: RPC server implementing all file operations.
    \item \textbf{client.py}: client application calling RPC functions.
    \item \textbf{server\_storage/}: server-side directory storing files.
\end{enumerate}

\section{Implementation}
Below is an example from the server:
\begin{lstlisting}[language=Python]
def upload_file(filename, data):
    filepath = os.path.join(BASE_DIR, filename)
    with open(filepath, "wb") as f:
        f.write(data.data)
    return True
\end{lstlisting}

\section{Group Work}
This practical was completed individually.

\section{Conclusion}
RPC simplifies networking by allowing the client to use remote functions as if they were local. The implementation successfully upgrades the TCP-based solution to an RPC-based file transfer system.

\end{document}
